\chapter{Paths and Files}

Exact paths follow \setting{data\_dir}, \setting{releases\_dir},
\setting{cache\_dir}, and XDG overrides. The table shows conceptual locations;
inspect the active manifest before recovery work.

\begin{shiptable}{L{0.36\textwidth}L{0.54\textwidth}}
\toprule
Path or file & Purpose \\
\midrule
\file{blackship.toml} & Default project manifest. \\
\file{Jailfile} / \file{Jailfile.toml} & Application image recipe. \\
\file{<data\_dir>/releases} & Directory-backed release roots or mountpoints. \\
\file{<data\_dir>/jails} & Jail roots or ZFS mountpoints. \\
\file{<data\_dir>/scope} & Per-jail ownership and rollback records. \\
\file{<data\_dir>/audit} & Durable per-jail JSONL histories. \\
\file{<data\_dir>/networks} & Named-network records. \\
\file{<data\_dir>/mcp.log} & MCP diagnostics when the server is used. \\
\file{<cache\_dir>} & Downloads and build cache material. \\
\file{<jail-root>/.blackship} & Staged jail definition and Jailfile carried by
snapshots/exports. \\
\file{/boot/loader.conf} & Boot-only RACCT enablement. \\
\file{/etc/pf.conf} & Host PF policy and Blackship anchor load points. \\
\file{/etc/rc.conf} & Host service enablement, including packaged warden setup. \\
\bottomrule
\end{shiptable}

\section*{Recovery rule}

Runtime state files are implementation records, not an alternate configuration
API. Back them up for diagnosis, but do not edit them to force desired state.
Prefer a validated manifest plus \command{eva}, \command{cleanup},
\command{setup}, or \command{gc}. If manual host repair is unavoidable, capture
the original files and audit trail first.
